%%%% Better Poster latex template example v1.0 (2019/04/04)
%%%% GNU General Public License v3.0
%%%% Rafael Bailo
%%%% https://github.com/rafaelbailo/betterposter-latex-template
%%%%
%%%% Original design from Mike Morrison
%%%% https://twitter.com/mikemorrison

\documentclass[a0paper,fleqn]{betterposter}
\usepackage{float}
\usepackage[latin1]{inputenc}
\usepackage{tikz}
\usepackage{booktabs}
\usetikzlibrary{shapes,arrows}
\usetikzlibrary{arrows.meta}

\usepackage{enumitem}
\usepackage{fontawesome}

\begin{document}

\betterposter{

  %%%%%%%% MAIN COLUMN

  \maincolumn{

    % Challenge

    {\fontsize{116}{150}\selectfont

      \faBullhorn~Even though teamwork is central to software engineering
      education, forming groups is difficult!

    }

    \vspace*{-1in}
    \par\noindent\rule{.25\textwidth}{8pt}
    \vspace*{1in}

    % Solution

    {\fontsize{116}{150}\selectfont

    \faLightbulbO~Ready for classroom adoption, \textbf{GatorGrouper} uses data
    to automatically form effective teams

    }

    \vspace*{-1in}
    \par\noindent\rule{.25\textwidth}{8pt}
    \vspace*{1.5in}

    % Equation

    {\fontsize{180}{80}\selectfont

      \faUser{} \, + \,
      \faDatabase{} \, + \,
      \faCode{} \, = \,
      \faUsers{}

    \vspace*{-5in}

    }

    }{

    %%%% Bottom space

    %% QR code
    %
    \qrcode{img/qr_code}{img/smartphoneWhite}{\textbf{Scan the QR Code} to
    \\visit our GitHub project!}

  }

  }{

  %%%%%%%% LEFT COLUMN

  \title{\fontsize{50}{60}\selectfont GatorGrouper: \\ \mbox{Using Python} and GitHub for
  Team Formation and Assessment}
  %
  \vspace*{-.25in}
  \author{Gregory M.\@ Kapfhammer \\[.5in]
  \mbox{\faLink{} gregorykapfhammer.com} \\
  \faTwitter{} @GregKapfhammer \\
  \faGithubAlt{} github.com/gkapfham
  %
  }

  \vspace*{-.5in}

  \section{\faSend~Introduction}
  %
  Python programs are often complex and difficult to test.
  %
  But test coverage does not show that a method was adequately tested.
  %
  A method $m$ is psuedo-tested if a test passes even when $m$ is not run.

  % \vspace*{.25in}

  % Since it may be time-consuming and error-prone to manually detect
  % psuedo-tested methods, Function-Fiasco automatically discovers them for
  % engineers.

  \vspace*{-.25in}
  %
  \section{\faCog~Algorithms}
  %
  GatorGrouper forms teams with:

  \begin{itemize}
    %
    \item Random
    \item Round-robin
    \item Kerninghan-Lin
    \item Genetic Algorithm
    %
  \end{itemize}

  %\vspace*{.5in}
  %%
  %Function-Fiasco performs these steps:
  %%
  %\begin{enumerate}[leftmargin=.5in]
  %  %
  %  \item Instrument all program methods
  %  \item Elide execution of chosen method
  %  \item Run the tests and observe behavior
  %  \item Run steps (2) and (3) for all methods
  %  \item Report the psuedo-tested methods
  %  %
  %\end{enumerate}

  \vspace*{.25in}
  %
  Steps are optimized, ensuring that the tool scales to large Python programs.
  %
  \vspace*{.5in}
  %

  %% This fills the space between the content and the logo
  \vfill

  \includegraphics[width=\textwidth]{img/CollegeLogo.eps}\\

  }{

  %%%%%%%% RIGHT COLUMN

  \vspace*{-.5in}
  %
  \section{\faCodeFork~Development}
  %
  We adopted these tools and processes:\\
  %
  \vspace*{-.5in}
  %
  \begin{itemize}[leftmargin=*]

    \item{Pytest for automated testing}
      %
    \item{Pyenv for language management}
      %
    \item{Pipenv for handling environments}
      %
    \item{Travis for integration and deployment}
      %
    \item{GitHub flow for project development}

  \end{itemize}

  \section{\faCloudUpload~Deployment}
  %
  You can use GatorGrouper these ways:\\
  %
  \vspace*{-.5in}
  %
  \begin{itemize}[leftmargin=*]

    \item{Command-line with local storage}
      %
    \item{Django-based web application}
      %
    \item{Leverage Amazon Elastic Beanstalk}

  \end{itemize}

  \section{\faDashboard~Future Work}
  %
  Add new features to GatorGrouper: \\
  %
  \vspace*{-.5in}

  \begin{itemize}[leftmargin=*]

    \item{Provide new ways to form teams}
      %
    \item{Accept and use feedback on teams}
      %
    \item{Leverage data from GatorGrader}
      %
    \item{Integrate with GitHub Classroom}

  \end{itemize}

  \vspace{.5em}
  %
  Considering both industry \& academia, use improved GatorGrouper to
  empirically study team effectiveness.

  \section{\faRocket~Action Steps}
  %
  Created with Python, GatorGrouper is a tool that supports team creation by
  balancing knowledge, skills, and conflicts.

  \vspace*{.25in}
  %
  If you would like support the development of GatorGrouper, please raise an
  issue on the tracker or create a pull request to add a new feature or bug fix.

  \vspace*{.25in}
  %
  See \faGithubAlt~\underline{GatorEducator/GatorGrader} for a related tool that
  automatically checks the work of writers and programmers.

  \section{\faThumbsOUp~Acknowledgements}
  %
  Check \faGithubAlt~\underline{GatorEducator/GatorGrouper} for a full list of
  project contributors.

  % \vspace*{.5in}
  % \includegraphics[width=\textwidth]{img/ComputerScience-Stack}

}

\end{document}
